\documentclass[12pt]{article}
\usepackage[utf8]{inputenc}
\usepackage[T1]{fontenc}
\usepackage{hyperref}
\title{Mili Language Resources for K-12 Curriculum Development}
\author{Compiled by Claude Sonnet 5.5 (claude-sonnet-5-5) for Andr\'as Kornai}
\date{}
\begin{document}
\maketitle

\section{Introduction}
Mili (ISO 639-3: ymh) is a Loloish language of the Sino-Tibetan family, spoken in Yunnan, China. Wikipedia places it in Jingdong, Yun, Zhenyuan and Xinping counties and describes it as a variety of Lolopo. Speaker counts vary: Wikipedia gives 23,000 (2002) and Joshua Project about 31,000. The Mili are officially counted within the Yi nationality. No educational materials in Mili were found.

\section{Primary Educational Resources}
No primers, textbooks or school programmes in Mili were found in the searches.

\section{Linguistic and Reference Materials}
\begin{itemize}
\item Wikipedia, Mili language: \url{https://en.wikipedia.org/wiki/Mili_language}
\item Wikipedia, Duota language (a related variety): \url{https://en.wikipedia.org/wiki/Duota_language}
\item Glottolog, Mili (mili1235): \url{https://glottolog.org/resource/languoid/id/mili1235.bigmap.html}
\end{itemize}

\section{Religious and Community Resources}
Joshua Project lists no Bible translation for Mili (status ``Translation Needed'').
\begin{itemize}
\item Joshua Project, Mili in China: \url{https://joshuaproject.net/people_groups/18596/CH}
\item Joshua Project, language page: \url{https://joshuaproject.net/languages/ymh}
\item Asia Harvest people-group profile of the Mili (PDF): \url{https://people-groups.asiaharvest.org/people-group-profiles/yunnan/mili.pdf}
\end{itemize}

\section{Digital and Audio Resources}
No audio, video or software resources specific to Mili were found.

\section{Curriculum Development Context}
Documentation consists of general reference pages only. No known school use, orthography or published teaching materials turned up. Standard Chinese and the Yi-nationality context would shape any programme. Curriculum work would need new documentation and community input.

\end{document}